% chapters/11/section-04.tex
\section{Menilai Mutu Perilaku}
\label{sec:menilai-mutu}

Inilah inti bab ini, yaitu menilai, bukan menulis. Mutu perilaku NPC dinilai
dari tiga segi, dan ketiganya diukur, bukan dikira-kira.

\subsection{Tiga Segi Penilaian}

\begin{description}
  \item[Kelancaran] Seberapa lancar permainannya berjalan. Ukurannya laju
        bingkai yang benar-benar tercapai. Apabila laju bingkai turun ketika
        jumlah NPC bertambah, beban perhitungannya sudah melampaui batas.
  \item[Daya tanggap] Seberapa cepat NPC menyesuaikan diri terhadap perpindahan
        pemain. Ukurannya jarak rata-rata NPC terhadap pemain selama permainan:
        jarak yang lebih kecil berarti NPC lebih rapat mendekati pemain.
  \item[Beban komputasi] Seberapa banyak pekerjaan yang diperlukan untuk
        menghasilkan perilaku itu. Ukurannya banyaknya pencarian jalur dan
        banyaknya petak yang diperiksa.
\end{description}

Ketiga segi itu saling bertentangan. Menambah daya tanggap berarti menambah
beban komputasi, dan beban yang berlebihan mengurangi kelancaran. Tugas
perancang adalah mencari titik keseimbangan, dan titik itu hanya dapat
ditemukan dengan mengukur.

\subsection{Hasil Pengukuran}

Program bab ini menyediakan pilihan \code{--evaluasi} yang menjalankan
permainan yang sama beberapa kali dengan selang hitung ulang yang berbeda, lalu
membandingkan hasilnya. Tabel~\ref{tab:evaluasi-selang} memuat hasilnya untuk
tiga NPC selama 900 bingkai.

\begin{table}[htbp]
\centering
\caption{Penilaian perilaku NPC pada beberapa selang hitung ulang jalur}
\label{tab:evaluasi-selang}
{\setlength{\tabcolsep}{4pt}\footnotesize%
\begin{tabularx}{\textwidth}{@{}C{2.0cm} C{2.4cm} C{2.6cm} Y@{}}
\toprule
\headrow \textbf{Selang} & \textbf{Hitung ulang} & \textbf{Petak diperiksa} & \textbf{Jarak rata-rata ke pemain} \\
\midrule
0,00 detik (setiap bingkai) & 2.700 & 251.463 & 35 piksel \\
\addlinespace
0,15 detik & 942 & 236.940 & 34 piksel \\
\addlinespace
0,30 detik & 858 & 240.582 & 34 piksel \\
\addlinespace
0,60 detik & 636 & 201.051 & 38 piksel \\
\addlinespace
1,20 detik & 594 & 189.237 & 40 piksel \\
\bottomrule
\end{tabularx}}
\tabsource{Pengukuran penulis: tiga NPC selama 900 bingkai dengan jalur pemain
yang tetap, sehingga seluruh baris dapat dibandingkan secara adil.}
\end{table}

\subsection{Menafsirkan Hasilnya}

Tiga hal dapat dibaca dari tabel tersebut, dan ketiganya merupakan hasil
penilaian yang sesungguhnya.

\begin{enumerate}
  \item \textbf{Pencarian setiap bingkai sangat boros.} Baris pertama
        menjalankan 2.700 pencarian, yaitu sekitar tiga kali lipat baris
        ketiga, tetapi jarak rata-ratanya hampir sama, yaitu 35 berbanding 34
        piksel. Pekerjaan tambahan sebesar tiga kali lipat itu tidak membawa
        perbaikan perilaku yang dapat dirasakan. Dengan kata lain,
        \textbf{sepertiga pekerjaan sudah cukup}.
  \item \textbf{Selang yang terlalu panjang mulai merugikan.} Pada selang 1,20
        detik, beban perhitungannya paling ringan, tetapi jarak rata-ratanya
        membesar menjadi 40 piksel. Artinya NPC mulai tertinggal, sebab
        jalurnya terlalu lama menuju tempat pemain berada sebelumnya.
  \item \textbf{Titik keseimbangannya berada di sekitar 0,30 detik.} Pada
        selang itu, beban perhitungannya tinggal sepertiga dibandingkan
        pencarian setiap bingkai, sedangkan jarak rata-ratanya justru yang
        terkecil. Inilah nilai yang dipakai sebagai bawaan pada program bab
        ini.
\end{enumerate}

\begin{peringatanbox}
Kesimpulan di atas berlaku untuk \textbf{peta, jumlah NPC, dan
jalur pemain yang diuji}, bukan untuk segala keadaan. Pada peta yang lebih
luas, pencariannya menjadi lebih mahal sehingga selang yang lebih panjang
mungkin lebih menguntungkan. Pada jalur pemain yang lebih berkelok, selang yang
lebih pendek mungkin diperlukan. Kesimpulan sebuah pengukuran selalu terikat
pada keadaan pengukurannya, dan menyebutkan keadaan itu merupakan bagian dari
laporan yang jujur.
\end{peringatanbox}
